\section{Limitations and next steps}
\label{sec:limitations}

The design of this thesis is deliberately controlled, it keeps the Calvano pricing environment fixed and varies only the information and timing frictions. That is the source of the main contribution, because it makes the interaction effects interpretable. It also defines the boundary of the claim. The thesis identifies mechanisms inside a stylised Q-learning environment; it does not provide a market-by-market quantitative forecast.

\begin{itemize}[topsep=2pt, itemsep=4pt]
    \item \textbf{Asymmetric frictions.} All three frictions are imposed symmetrically on the two agents. Real markets often feature asymmetric frictions, for example one firm with a faster price-update cadence, cleaner data feed, or stronger scraping infrastructure than the other. \citet{brown2023}'s empirical findings on asymmetric update frequencies make this the most natural next extension.
    \item \textbf{Richer market structure.} The model uses a symmetric two-firm environment to isolate the information and timing mechanisms. The next step is to test whether the same interaction patterns survive with more firms, asymmetric costs, product-level heterogeneity, capacity constraints, or consumer search frictions.
    \item \textbf{Learning technology and memory.} Appendix~\ref{app:bounded-memory} explains why extending tabular Q-learning from one-period memory to two-period memory is computationally demanding in this state space. A meaningful extension would require either substantially larger compute, for example a high-performance-computing allocation, or a methodological switch to function-approximated $Q$. That step would also allow comparison with more modern pricing algorithms, but at the cost of giving up the transparent tabular mechanism used here.
\end{itemize}

These limitations are best read as next steps rather than qualifications of the core result. The contribution of the thesis is to show, in a controlled environment, that information and timing frictions cannot be evaluated one at a time, because their combination can qualitatively change the degree of learned collusion. Future work should test how far that mechanism travels once firm asymmetry, richer products, and more flexible learning rules are added back in.
